\documentclass[11pt,a4paper]{article}

\usepackage[main=spanish]{babel}
\usepackage[utf8]{inputenc}
\usepackage[T1]{fontenc}

\usepackage{amsmath, amssymb, amsfonts}
\usepackage{mathtools}
\usepackage{geometry}
\usepackage{graphicx}
\usepackage{setspace}
\usepackage{hyperref}
\usepackage{graphicx}
\usepackage{appendix}
\usepackage{listings}
\usepackage{xcolor}

\geometry{margin=2.5cm}
\onehalfspacing
\title{Laboratorio 1: Optimización de la función de Rosenbrock}
\author{Adrián Janus Gonzalez Adamuz}
\date{\today}
\lstset{
  inputencoding=utf8,
  extendedchars=true,
  basicstyle=\ttfamily\small,
  numbers=left,
  numberstyle=\tiny,
  frame=single,
  breaklines=true,
  showstringspaces=false,
  tabsize=2,
  literate=
    {á}{{\'a}}1 {é}{{\'e}}1 {í}{{\'\i}}1 {ó}{{\'o}}1 {ú}{{\'u}}1
    {Á}{{\'A}}1 {É}{{\'E}}1 {Í}{{\'I}}1 {Ó}{{\'O}}1 {Ú}{{\'U}}1
    {ñ}{{\~n}}1 {Ñ}{{\~N}}1
  }	

\begin{document}
\maketitle

\section{La función de Rosenbrock}

La función de Rosenbrock: 
\[
f:\mathbb{R}^2 \to \mathbb{R}
\]
definida por
\begin{equation}
f(x,y) = (a - x)^2 + b\,(y - x^2)^2,
\qquad a>0,\; b>0.
\label{eq:rosenbrock}
\end{equation}

Dicha función es no negativa para todo $(x,y)\in\mathbb{R}^2$, pues se expresa como suma de cuadrados. La geometría de su gráfica está dominada por un valle curvo que se aproxima a la parábola $y=x^2$. 

\subsection{Puntos estacionarios}

Para determinar los puntos estacionarios de la función \eqref{eq:rosenbrock}, se calcula su gradiente. Las derivadas parciales vienen dadas por
\[
\frac{\partial f}{\partial x}(x,y) = -2(a-x) - 4bx(y-x^2),
\qquad
\frac{\partial f}{\partial y}(x,y) = 2b(y-x^2).
\]

Un punto $(x,y)$ es estacionario si satisface $\nabla f(x,y)=0$. De la ecuación $\partial f/\partial y=0$ se obtiene 
\[
y = x^2.
\]
Al sustituir en $\partial f/\partial x=0$, se obtiene
\[
-2(a-x)=0,
\]
lo cual implica $x=a$. En consecuencia, el único punto estacionario de la función de Rosenbrock es
\[
(x^\ast,y^\ast) = (a,a^2).
\]

\subsection{Clasificación del punto estacionario}

Para clasificar la naturaleza del punto estacionario encontrado, se considera la matriz Hessiana de $f$, cuyas entradas están dadas por
\[
\frac{\partial^2 f}{\partial x^2} = 2 - 4by + 12bx^2,
\qquad
\frac{\partial^2 f}{\partial x \partial y} = -4bx,
\qquad
\frac{\partial^2 f}{\partial y^2} = 2b.
\]

Evaluando la Hessiana en el punto $(a,a^2)$ se obtiene
\[
Hf(a,a^2) =
\begin{pmatrix}
2 + 8ba^2 & -4ba \\
-4ba & 2b
\end{pmatrix}.
\]

El determinante de esta matriz es
\[
\det Hf(a,a^2) = (2+8ba^2)(2b) - (4ba)^2 = 4b > 0,
\]
y además se cumple que
\[
\frac{\partial^2 f}{\partial x^2}(a,a^2) = 2 + 8ba^2 > 0.
\]
El Hessiano es positivo definido, por lo que el punto $(a,a^2)$ es un mínimo local estricto. Dado que $f(x,y)\geq 0$ para todo $(x,y)$ y que $f(a,a^2)=0$, dicho mínimo es además un mínimo global.

\section{Problema de optimización con restricción}

Se considera ahora el mismo problema de optimización, pero sujeto a la restricción de igualdad
\begin{equation}
x^2 + y^2 = c^2,
\qquad c>0,
\label{eq:restriccion}
\end{equation}
la cual describe una circunferencia de radio $c$ centrada en el origen.

\subsection{Lagrangiano}

Para abordar el problema de optimización restringida se introduce el Lagrangiano
\[
\mathcal{L}(x,y,\lambda)
= (a-x)^2 + b(y-x^2)^2 + \lambda(x^2+y^2-c^2),
\]


Las condiciones necesarias de primer orden vienen dadas por el sistema
\[
\frac{\partial \mathcal{L}}{\partial x} = -2(a-x) - 4bx(y-x^2) + 2\lambda x = 0,
\]
\[
\frac{\partial \mathcal{L}}{\partial y} = 2b(y-x^2) + 2\lambda y = 0,
\]
\[
\frac{\partial \mathcal{L}}{\partial \lambda} = x^2 + y^2 - c^2 = 0.
\]

Este sistema no lineal de tres ecuaciones en las incógnitas $(x,y,\lambda)$ no admite, en general, una solución analítica cerrada, por lo que su resolución se lleva a cabo mediante métodos numéricos.

\subsection{Resolución numérica}
Al variar el parámetro $c$, se obtiene una familia de puntos estacionarios que describe una trayectoria en el plano $(x,y)$, la cual refleja cómo se desplaza el óptimo restringido conforme cambia la geometría de la restricción.
\begin{figure}[hbtp]
\centering
\includegraphics[scale=0.3]{../../Imágenes/Graficas/Rosenbrock.png} 
\caption{Resultados Completos}
\end{figure}

\newpage
\section{Hiperbanana}
La función de generalizada de Rosenbrock: 
\[
f:\mathbb{R}^n \to \mathbb{R}^n
\]
definida 
\[
f(x) = \sum_{i=1}^{N-1} \bigl( (a - x_i)^2 + b (x_{i+1} - x_i^2)^2 \bigr), a>0, b>0
\]
para $a=1$ y $b=100$

\subsection{Caso N=3}
Para \(N=3\), con \(a=1\) y \(b=100\), la función toma la forma
\begin{equation}
f(x_1,x_2,x_3)
=(1-x_1)^2+100(x_2-x_1^2)^2
+(1-x_2)^2+100(x_3-x_2^2)^2.
\end{equation}
Dado que \(f:\mathbb{R}^3\to\mathbb{R}\), la visualización directa mediante una superficie no es posible. En su lugar, se estudia su comportamiento geométrico mediante cortes y proyecciones, fijando una variable y analizando curvas de nivel en los planos restantes, así como evaluando la función sobre trayectorias relevantes.


\begin{figure}[hbtp]
\centering
\includegraphics[scale=0.3]{../../Imágenes/Rosen_gen.png} 
\caption{Resultados Completos}
\end{figure}





\subsection{Condiciones de primer orden para \(N\) general}

Los puntos estacionarios de \(f\) satisfacen el sistema no lineal \(\nabla f(x)=0\). Las derivadas parciales están dadas por:
\begin{align}
\frac{\partial f}{\partial x_1}
&=-2(a-x_1)-4b\,x_1(x_2-x_1^2),\\
\frac{\partial f}{\partial x_i}
&=-2(a-x_i)-4b\,x_i(x_{i+1}-x_i^2)+2b(x_i-x_{i-1}^2),
\quad 2\le i\le N-1,\\
\frac{\partial f}{\partial x_N}
&=2b(x_N-x_{N-1}^2).
\end{align}
Por lo tanto, el sistema de ecuaciones para un punto estacionario se escribe como
\begin{equation}
\nabla f(x)=0.
\end{equation}
Una posible solución es:
\[
x_1=a,\qquad x_{i+1}=x_i^2,\quad i=1,\dots,N-1.
\]
En el caso particular \(a=1\), esta relación implica
\[
x^\star=(1,1,\dots,1),
\]
el cual corresponde al mínimo global de la función, con valor \(f(x^\star)=0\).

\subsection{Resolución numérica del sistema sin restricción}

Para \(a=1\), \(b=100\) y \(N=100\), el sistema \(\nabla f(x)=0\) es un sistema no lineal de gran dimensión que no admite solución analítica explícita. Para resolverlo numéricamente se utilizó la función \texttt{root} del módulo \texttt{scipy.optimize}, la cual implementa distintos métodos para resolver sistemas no lineales.

En particular, el problema se formula como la búsqueda de una raíz de la función
\[
F(x)=\nabla f(x),
\]
y se resuelve mediante un método tipo Newton o cuasi-Newton, según la configuración elegida internamente por \texttt{root}. El criterio de paro se basa en la norma del residuo \(\|F(x)\|\) y en la estabilidad de las iteraciones sucesivas.

Para los puntos dados conseguimos el siguiente resultado:
\begin{verbatim}
Converge: False
Norma del residuo ||grad||: 2.773162768811204
Primeras componentes:  [0.99999675 0.99999487 0.99999281 0.99999073 0.99998886]
\end{verbatim}
\section{Problema con restricción esférica}

Se considera ahora el mismo problema de optimización sujeto a la restricción de igualdad
\begin{equation}
\|x\|^2=c^2.
\end{equation}
Geométricamente, esta restricción define una esfera de radio \(c\) centrada en el origen en \(\mathbb{R}^N\).

\subsection{Lagrangiano y condiciones de primer orden}

El Lagrangiano asociado al problema es
\begin{equation}
\mathcal{L}(x,\lambda)
=f(x)+\lambda\left(\sum_{i=1}^N x_i^2-c^2\right).
\end{equation}
\begin{align}
\nabla_x \mathcal{L}(x,\lambda)
&=\nabla f(x)+2\lambda x=0,\\
\sum_{i=1}^N x_i^2&=c^2.
\end{align}

Al escribir las ecuaciones componente a componente se obtiene un sistema no lineal de \(N+1\) ecuaciones. Dependiendo del valor del multiplicador de Lagrange \(\lambda\) y de la posible presencia de componentes nulas en \(x\), el sistema presenta tres casos cualitativamente distintos, los cuales deben tratarse con cuidado en el análisis numérico.

\subsection{Resolución numérica del sistema KKT}

El sistema de condiciones de primer orden se resolvió numéricamente utilizando nuevamente la función \texttt{root} de \texttt{scipy.optimize}, aplicada ahora al sistema aumentado
\[
F(x,\lambda)=
\begin{pmatrix}
\nabla f(x)+2\lambda x\\
\sum_{i=1}^N x_i^2-c^2
\end{pmatrix}.
\]
Este enfoque corresponde a aplicar un método de Newton al sistema KKT del problema restringido. La convergencia depende de manera sensible de la elección del punto inicial, tanto para \(x\) como para \(\lambda\), así como del valor del parámetro \(c\).
\begin{verbatim}

Converge: False
||grad f + 2lambdax||: 1.3006886874590542
||x||^2: 4.085219067280299   c^2: 4.0

\end{verbatim}

\subsection{Barrido en el parámetro \(c\)}
El procedimiento anterior se repitió para distintos valores del parámetro \(c\), con \(c\in\{0,0.1,0.2,\dots,5\}\). Para cada valor de \(c\) se obtuvo un punto estacionario factible \(x^\star(c)\).

Dado que la restricción impone \(\|x^\star(c)\|=c\), la distancia del mínimo al origen coincide necesariamente con el valor de \(c\). Por esta razón, además de verificar esta relación numéricamente, resulta más informativo analizar el comportamiento del valor óptimo \(f(x^\star(c))\) o la distancia \(\|x^\star(c)-\mathbf{1}\|\), donde \(\mathbf{1}=(1,\dots,1)\) es el minimizador no restringido. Estos resultados permiten cuantificar cómo la restricción geométrica afecta la localización y el valor del mínimo.
\begin{verbatim}
Total fallas: 9
Último c exitoso: 4.4

Resultados por valor de c:
c =  0.0  ->  Éxito
c =  0.1  ->  Éxito
c =  0.2  ->  Éxito
c =  0.3  ->  Éxito
c =  0.4  ->  Éxito
c =  0.5  ->  Éxito
c =  0.6  ->  Falla
c =  0.7  ->  Falla
c =  0.8  ->  Falla
c =  0.9  ->  Éxito
c =  1.0  ->  Éxito
c =  1.1  ->  Éxito
c =  1.2  ->  Éxito
c =  1.3  ->  Éxito
c =  1.4  ->  Éxito
c =  1.5  ->  Éxito
c =  1.6  ->  Éxito
c =  1.7  ->  Éxito
c =  1.8  ->  Éxito
c =  1.9  ->  Éxito
c =  2.0  ->  Éxito
c =  2.1  ->  Éxito
c =  2.2  ->  Éxito
c =  2.3  ->  Éxito
c =  2.4  ->  Éxito
c =  2.5  ->  Éxito
c =  2.6  ->  Éxito
c =  2.7  ->  Éxito
c =  2.8  ->  Éxito
c =  2.9  ->  Éxito
c =  3.0  ->  Éxito
c =  3.1  ->  Éxito
c =  3.2  ->  Éxito
c =  3.3  ->  Éxito
c =  3.4  ->  Éxito
c =  3.5  ->  Éxito
c =  3.6  ->  Éxito
c =  3.7  ->  Éxito
c =  3.8  ->  Éxito
c =  3.9  ->  Éxito
c =  4.0  ->  Éxito
c =  4.1  ->  Éxito
c =  4.2  ->  Éxito
c =  4.3  ->  Éxito
c =  4.4  ->  Éxito
c =  4.5  ->  Falla
c =  4.6  ->  Falla
c =  4.7  ->  Falla
c =  4.8  ->  Falla
c =  4.9  ->  Falla
c =  5.0  ->  Falla
\end{verbatim}

\appendix

\section{Código para visualización y problema no restringido}
\label{ap:codigo1}

\lstinputlisting{laboratorio1.py}
\section{Código para el problema restringido y método de continuación}
\label{ap:codigo2}

\lstinputlisting{laboratorio1_1.py}

\end{document}
